% Titulos de tablas y figuras (espanol). Compartidos por los documentos y por las versiones PNG.
\newcommand{\capTabA}{Supervivencia de plántulas en LHS-1 según la proporción de sustrato ionitopónico. Los valores se expresan como media $\pm$ desviación estándar (DE); $n = 5$ unidades experimentales por grupo.}
\newcommand{\capTabB}{Análisis de varianza para la proporción de sustrato ionitopónico en LHS-1, particionado en dos contrastes ortogonales planificados.}
\newcommand{\capTabC}{Contrastes planificados para la supervivencia de plántulas en LHS-1. Las diferencias se expresan en puntos porcentuales (pp).}
\newcommand{\capTabD}{Sensibilidad de los contrastes planificados a la elección del término de error y de la unidad de análisis.}
\newcommand{\capTabE}{Contrastes de contexto para la supervivencia de plántulas. Las diferencias se expresan en puntos porcentuales (pp).}
\newcommand{\capTabF}{Análisis de varianza de dos vías para la supervivencia de plántulas según sustrato y tratamiento.}
\newcommand{\capTabG}{Germinación según tratamiento y sustrato. Los valores se expresan como proporciones con intervalos de confianza (IC) exactos (Clopper--Pearson) del 95\,\%.}
\newcommand{\capTabH}{Contrastes seleccionados para la germinación, evaluados a nivel de semilla y a nivel de unidad experimental.}
\newcommand{\capTabI}{Potencia estadística de la fase piloto y tamaño de muestra requerido para un experimento confirmatorio (potencia de 0.80, $\alpha = 0.05$).}
\newcommand{\capFigA}{Supervivencia de plántulas en LHS-1 según la proporción de sustrato ionitopónico. Los puntos representan las medias de grupo y las barras representan $\pm 1$ error estándar (error agrupado, 32 grados de libertad); $n = 5$ unidades experimentales por grupo.}
\newcommand{\capFigB}{Diferencias estimadas en la supervivencia de plántulas con intervalos de confianza bilaterales del 95\,\%. Los marcadores oscuros corresponden a los contrastes planificados de la hipótesis; los marcadores grises corresponden a comparaciones complementarias y de contexto.}
\newcommand{\capFigC}{Supervivencia de plántulas según tratamiento y sustrato. Los puntos representan las medias de grupo y las barras representan $\pm 1$ error estándar (error agrupado, 32 grados de libertad).}
\newcommand{\capFigD}{Germinación según tratamiento y sustrato. Los puntos representan las proporciones observadas y las barras representan intervalos de confianza exactos (Clopper--Pearson) del 95\,\%; 20 semillas por grupo.}
\newcommand{\capFigE}{Potencia estadística para detectar una diferencia de 20 puntos porcentuales en la supervivencia de plántulas en función del número de unidades experimentales por grupo ($\alpha = 0.05$; desviación estándar residual de 31.8 puntos porcentuales). La línea horizontal indica una potencia de 0.80.}
\newcommand{\capFigTA}{Distribución $t$ de Student con 32 grados de libertad bajo la hipótesis nula de C1. El estadístico observado no alcanza la región de rechazo unilateral.}
\newcommand{\capFigTB}{Distribución de permutación exacta de la diferencia de germinación entre el aditivo al 5\,\% y la solución convencional en LHS-1, con la unidad experimental como unidad de análisis (252 reordenamientos de 10 unidades). En azul, los reordenamientos al menos tan extremos como el observado.}
